\scanpage{31}
\mspage{5/}

canonique stable par translations etc, mais une classe de telles mod homothétie par $c\in\mathbb R^{*+}$ -- ces métriques sont caractérisées par la condition d'être invariantes par translation [\uline{NB} $X$ a automatiquement une structure de droite affine sur $\mathbb C$]. Ces métriques \deleted{sont \unsure{caractérisées}} ne dépendent que de la structure conforme sous-jacente.

3) \uline{Cas} \deleted{\unsure{de la sphère} $X\simeq\mathbb P^1_{\mathbb C}$} $X\simeq$ \deleted{\unsure{droite|sphère}}plan projecti\deleted{ve}f réel. Dans ce \deleted{premier} \deleted{\illegible{1}} cas, $X^{\mathrm{or}}$ est une droite projective complexe munie d'un antiautomorphisme invol. $\sigma$ \uline{sans pt fixe} (antipodisme), ou ce qui revient au $=$, d'une métrique ``euclidienne sphérique'' compatible avec sa structure conforme. [Celle-ci, \deleted{étant stable} \inserted{\unsure{supp.|sera} normalisée} \deleted{par l'antipodisme,} \inserted{par la condition} que l'aire de la surface soit $=1$] É\deleted{\unsure{tant}}tant invariante par l'antipodisme, cette métrique se descend à $X$.

$3')$ Cas $X\simeq\mathbb P^1_{\mathbb C}$. Il n'y a pas de métrique riemannienne sur $X$, invariante [ou quasi-invariante] par tt automorphisme -- \unsure{dem.|demande|donc} une structure suppl. sur $X$

% archive sheet number: 33
